\chapter{Boolean Analysis — Zk Fourier}
%%% generated by the blueprint pipeline from TCSlib.BooleanAnalysis.BLR.ZkFourier
\section{Overview}

This module develops the Fourier analysis of complex-valued functions on the
domain $\mathbb{Z}_k^n = (\mathbb{Z}/k\mathbb{Z})^n$, the natural
generalization of the Boolean hypercube $\{0,1\}^n$.  It establishes the
primitive $k$-th root of unity $\omega_k = e^{2\pi i/k}$, the associated
group homomorphism $a \mapsto \omega_k^a$, the Fourier characters
$\chi_s(x) = \omega_k^{s \cdot x}$, and the core identities: orthonormality
of characters, the Fourier expansion, Parseval's identity, and the convolution
theorem.

%%%%%%%%%%%%%%%%%%%%%%%%%%%%%%%%%%%%%%%%%%%%%%%%%%%%%%%%%%%%%%%%%%%%%%%%%%%%%%%
\section{Declarations}

\begin{definition}[Vectors over $\mathbb{Z}/k\mathbb{Z}$]
\lean{ZkFourier.ZkVec}
\label{ZkFourier.ZkVec}
\leanok
$\texttt{ZkFourier.ZkVec}\ k\ n$ is the $n$-dimensional vector space over
$\mathbb{Z}/k\mathbb{Z}$, defined as the function type
$\mathrm{Fin}\,n \to \mathbb{Z}/k\mathbb{Z}$.  It carries the pointwise
additive group structure of $(\mathbb{Z}/k\mathbb{Z})^n$.
\end{definition}

\begin{lemma}[Cardinality of the space of vectors over $\mathbb{Z}/k\mathbb{Z}$]
\lean{ZkFourier.card_ZkVec}
\label{ZkFourier.card_ZkVec}
\leanok
\uses{ZkFourier.ZkVec, ZkFourier.rootOfUnity}
\difficulty{1}
Let $k$ and $n$ be natural numbers with $k \ge 1$. The set $(\mathbb{Z}/k\mathbb{Z})^n$
of $n$-tuples with entries in $\mathbb{Z}/k\mathbb{Z}$ is finite, and its cardinality is
\[
\bigl|(\mathbb{Z}/k\mathbb{Z})^n\bigr| = k^n.
\]
\end{lemma}

\begin{definition}[Primitive $k$-th root of unity]
\lean{ZkFourier.rootOfUnity}
\label{ZkFourier.rootOfUnity}
\leanok
\statementsource{boolean-ch08-generalized-domains}{pdf-50833afe7403-p022-b007}
$\texttt{ZkFourier.rootOfUnity}\ k$ is the complex number
$\omega_k = e^{2\pi i / k}$, the primitive $k$-th root of unity used as the
base of the Fourier characters on $\mathbb{Z}_k^n$.
\end{definition}

\begin{definition}[Embedding of $\mathbb{Z}/k\mathbb{Z}$ into $\mathbb{C}^\times$]
\lean{ZkFourier.toOmega}
\label{ZkFourier.toOmega}
\leanok
\statementsource{boolean-ch08-generalized-domains}{pdf-50833afe7403-p022-b007}
\uses{ZkFourier.rootOfUnity}
For $a : \mathbb{Z}/k\mathbb{Z}$, $\texttt{ZkFourier.toOmega}\ a = \omega_k^{a}$,
where $a$ is lifted to its canonical representative in
$\{0, \ldots, k-1\} \subset \mathbb{N}$.  This is the group homomorphism
$(\mathbb{Z}/k\mathbb{Z}, +) \to (\mathbb{C}^\times, \cdot)$.
\end{definition}

\begin{lemma}[$\omega_k$ is a primitive $k$-th root of unity]
\lean{ZkFourier.isPrimitiveRoot_rootOfUnity}
\label{ZkFourier.isPrimitiveRoot_rootOfUnity}
\leanok
\uses{ZkFourier.rootOfUnity}
\difficulty{1}
Let $k$ be a positive integer and set $\omega_k = e^{2\pi i / k} \in \bbc$. Then
$\omega_k$ is a primitive $k$-th root of unity: $\omega_k^k = 1$, and no smaller
positive power of $\omega_k$ equals $1$, so that $\omega_k^j = 1$ for a natural number
$j$ holds precisely when $k$ divides $j$.
\end{lemma}

\begin{lemma}[The $k$-th power of a $k$-th root of unity]
\lean{ZkFourier.rootOfUnity_pow_k}
\label{ZkFourier.rootOfUnity_pow_k}
\leanok
\uses{ZkFourier.isPrimitiveRoot_rootOfUnity, ZkFourier.rootOfUnity}
\difficulty{1}
Let $k$ be a positive integer and let $\omega_k = e^{2\pi i / k}$ be the primitive
$k$-th root of unity in $\mathbb{C}$. Then $\omega_k^k = 1$.
\end{lemma}

\begin{lemma}[The additive identity maps to one]
\lean{ZkFourier.toOmega_zero}
\label{ZkFourier.toOmega_zero}
\leanok
\uses{ZkFourier.toOmega}
\difficulty{1}
Let $k \ge 1$ and write $\omega_k = e^{2\pi i/k}$ for the primitive $k$-th root of
unity, and for $a \in \mathbb{Z}/k\mathbb{Z}$ let its image under the homomorphism
$(\mathbb{Z}/k\mathbb{Z}, +) \to (\mathbb{C}^\times, \cdot)$ be $\omega_k^{a}$. Then the
additive identity $0 \in \mathbb{Z}/k\mathbb{Z}$ maps to $1$; that is, $\omega_k^{0} =
1$.
\end{lemma}

\begin{lemma}[Additivity of the root-of-unity embedding]
\lean{ZkFourier.toOmega_add}
\label{ZkFourier.toOmega_add}
\leanok
\uses{ZkFourier.rootOfUnity_pow_k, ZkFourier.toOmega}
\difficulty{3}
Fix an integer $k \ge 1$ and let $\omega_k = e^{2\pi i/k}$ be the primitive $k$-th root
of unity. For $a \in \mathbb{Z}/k\mathbb{Z}$ write $\chi(a) = \omega_k^{a}$, where the
exponent is the canonical representative of $a$ in $\{0, 1, \ldots, k-1\}$. Then for all
$a, b \in \mathbb{Z}/k\mathbb{Z}$,
\[
\chi(a + b) = \chi(a)\,\chi(b),
\]
so $\chi$ is a homomorphism from the additive group $\mathbb{Z}/k\mathbb{Z}$ to the
multiplicative group $\mathbb{C}^{\times}$.
\end{lemma}

\begin{lemma}[Conjugation of a Fourier character under negation]
\lean{ZkFourier.toOmega_neg}
\label{ZkFourier.toOmega_neg}
\leanok
\uses{ZkFourier.rootOfUnity, ZkFourier.toOmega}
\difficulty{5}
Let $k \ge 1$ and let $\omega_k = e^{2\pi i/k}$ be the primitive $k$-th root of unity,
and for $a \in \mathbb{Z}/k\mathbb{Z}$ write $\omega_k^{a}$ for the value $\omega_k$
raised to the canonical representative of $a$ in $\{0,\dots,k-1\}$. Then for every $a
\in \mathbb{Z}/k\mathbb{Z}$,
\[
\omega_k^{-a} = \overline{\omega_k^{a}},
\]
that is, replacing $a$ by its additive inverse conjugates the corresponding character
value.
\end{lemma}

\begin{lemma}[Unit modulus of powers of a root of unity]
\lean{ZkFourier.norm_toOmega}
\label{ZkFourier.norm_toOmega}
\leanok
\uses{ZkFourier.rootOfUnity, ZkFourier.toOmega}
\difficulty{1}
Fix an integer $k \ge 1$ and write $\omega_k = e^{2\pi i/k}$ for the primitive $k$-th
root of unity. For every residue $a \in \mathbb{Z}/k\mathbb{Z}$, the complex number
$\omega_k^{a}$, obtained by raising $\omega_k$ to the canonical representative of $a$ in
$\{0,\dots,k-1\}$, satisfies $\norm{\omega_k^{a}} = 1$, so it lies on the unit circle.
\end{lemma}

\begin{lemma}[Multiplicativity of the root-of-unity embedding]
\lean{ZkFourier.toOmega_mul}
\label{ZkFourier.toOmega_mul}
\leanok
\uses{ZkFourier.rootOfUnity, ZkFourier.rootOfUnity_pow_k, ZkFourier.toOmega}
\difficulty{3}
Fix a positive integer $k$ and let $\omega_k = e^{2\pi i / k}$. For $b \in
\mathbb{Z}/k\mathbb{Z}$, write $\lfloor b\rfloor \in \{0,\ldots,k-1\}$ for its canonical
representative and set $\chi(b) = \omega_k^{\lfloor b\rfloor}$. Then for all $j, a \in
\mathbb{Z}/k\mathbb{Z}$,
\[
  \chi(j \cdot a) = \bigl(\omega_k^{\lfloor j\rfloor}\bigr)^{\lfloor a\rfloor}.
\]
\end{lemma}

\begin{lemma}[Character orthogonality on $\mathbb{Z}/k\mathbb{Z}$]
\lean{ZkFourier.geom_sum_toOmega}
\statementsource{boolean-ch08-generalized-domains}{
  pdf-50833afe7403-p022-b002,
  pdf-50833afe7403-p022-b003
}
\label{ZkFourier.geom_sum_toOmega}
\leanok
\uses{ZkFourier.isPrimitiveRoot_rootOfUnity, ZkFourier.rootOfUnity, ZkFourier.rootOfUnity_pow_k, ZkFourier.toOmega, ZkFourier.toOmega_mul}
\difficulty{5}
Let $k \ge 1$ and let $\omega_k = e^{2\pi i/k}$ be the primitive $k$-th root of unity.
For each $a \in \mathbb{Z}/k\mathbb{Z}$, write $\omega_k^{a}$ for $\omega_k$ raised to
the canonical representative of $a$ in $\{0,\ldots,k-1\}$. Then for every $j \in
\mathbb{Z}/k\mathbb{Z}$,
\[
  \sum_{a \in \mathbb{Z}/k\mathbb{Z}} \omega_k^{\,j a}
  = \begin{cases} k & \text{if } j = 0, \\ 0 & \text{if } j \neq 0. \end{cases}
\]
\end{lemma}

\begin{definition}[Functions on $\mathbb{Z}_k^n$]
\lean{ZkFourier.ZkFun}
\label{ZkFourier.ZkFun}
\leanok
\statementsource{boolean-ch08-generalized-domains}{pdf-50833afe7403-p021-b002}
\uses{ZkFourier.ZkVec}
$\texttt{ZkFourier.ZkFun}\ k\ n$ is the type of complex-valued functions on
$\mathbb{Z}_k^n$, defined as $\mathbb{Z}_k^n \to \mathbb{C}$.  These are the
objects whose Fourier transforms are studied.
\end{definition}

\begin{definition}[Expectation]
\lean{ZkFourier.expectation}
\label{ZkFourier.expectation}
\leanok
\uses{ZkFourier.ZkFun, ZkFourier.ZkVec}
The expectation of $f : \mathbb{Z}_k^n \to \mathbb{C}$ is the uniform average
\[
  \mathbb{E}[f] = \frac{1}{k^n} \sum_{x \in \mathbb{Z}_k^n} f(x).
\]
\end{definition}

\begin{definition}[Hermitian inner product]
\lean{ZkFourier.inner_product}
\label{ZkFourier.inner_product}
\leanok
\statementsource{boolean-ch08-generalized-domains}{pdf-50833afe7403-p021-b002}
\uses{ZkFourier.ZkFun, ZkFourier.expectation}
The Hermitian inner product of $f, g : \mathbb{Z}_k^n \to \mathbb{C}$ is
\[
  \langle f, g \rangle = \mathbb{E}\!\left[f(x)\,\overline{g(x)}\right].
\]
\end{definition}

\begin{definition}[Squared $L^2$ norm]
\lean{ZkFourier.L2_norm_sq}
\label{ZkFourier.L2_norm_sq}
\leanok
\uses{ZkFourier.ZkFun, ZkFourier.ZkVec}
The squared $L^2$ norm of $f : \mathbb{Z}_k^n \to \mathbb{C}$ is the
real-valued quantity
\[
  \|f\|_2^2 = \frac{1}{k^n} \sum_{x \in \mathbb{Z}_k^n} |f(x)|^2.
\]
\end{definition}

\begin{definition}[Convolution]
\lean{ZkFourier.convolution}
\label{ZkFourier.convolution}
\leanok
\statementsource{boolean-ch08-generalized-domains}{pdf-50833afe7403-p023-b004}
\uses{ZkFourier.ZkFun, ZkFourier.expectation}
The convolution of $f, g : \mathbb{Z}_k^n \to \mathbb{C}$ is
\[
  (f * g)(x) = \mathbb{E}_{y}\!\left[f(y)\,g(x - y)\right].
\]
\end{definition}

\begin{definition}[Dot product on $\mathbb{Z}_k^n$]
\lean{ZkFourier.zkDot}
\label{ZkFourier.zkDot}
\leanok
\uses{ZkFourier.ZkVec}
For $s, x \in \mathbb{Z}_k^n$, the dot product is
$s \cdot x = \sum_{i=0}^{n-1} s_i \, x_i \in \mathbb{Z}/k\mathbb{Z}$,
computed componentwise modulo $k$.
\end{definition}

\begin{definition}[Fourier character]
\lean{ZkFourier.char_s}
\label{ZkFourier.char_s}
\leanok
\statementsource{boolean-ch08-generalized-domains}{
  pdf-50833afe7403-p022-b007,
  pdf-50833afe7403-p022-b009
}
\uses{ZkFourier.ZkFun, ZkFourier.ZkVec, ZkFourier.toOmega, ZkFourier.zkDot}
For $s \in \mathbb{Z}_k^n$, the Fourier character indexed by $s$ is
$\chi_s : \mathbb{Z}_k^n \to \mathbb{C}$ defined by
$\chi_s(x) = \omega_k^{s \cdot x}$.
The $k^n$ characters $\{\chi_s\}_{s \in \mathbb{Z}_k^n}$ form an orthonormal
basis of $L^2(\mathbb{Z}_k^n)$.
\end{definition}

\begin{lemma}[Right-additivity of the dot product over $\mathbb{Z}/k\mathbb{Z}$]
\lean{ZkFourier.zkDot_add_right}
\label{ZkFourier.zkDot_add_right}
\leanok
\uses{ZkFourier.ZkVec, ZkFourier.zkDot}
\difficulty{2}
Fix an integer $k \ge 1$ and a dimension $n$, and let $s, x, y \in
(\mathbb{Z}/k\mathbb{Z})^n$. Writing $s \cdot x = \sum_{i=0}^{n-1} s_i x_i \in
\mathbb{Z}/k\mathbb{Z}$ for the componentwise dot product, the dot product is additive
in its right argument:
\[
s \cdot (x + y) = s \cdot x + s \cdot y,
\]
where $x + y$ denotes the componentwise sum.
\end{lemma}

\begin{lemma}[Dot product with the zero vector]
\lean{ZkFourier.zkDot_zero_left}
\label{ZkFourier.zkDot_zero_left}
\leanok
\uses{ZkFourier.ZkVec, ZkFourier.zkDot}
\difficulty{1}
Fix an integer $k \ge 1$ and let $x \in (\bbz/k\bbz)^n$ be any vector. Then the dot
product of the zero vector with $x$ is zero: $0 \cdot x = 0$ in $\bbz/k\bbz$.
\end{lemma}

\begin{lemma}[Dot product against the zero vector]
\lean{ZkFourier.zkDot_zero_right}
\label{ZkFourier.zkDot_zero_right}
\leanok
\uses{ZkFourier.ZkVec, ZkFourier.zkDot}
\difficulty{1}
Fix an integer $k \ge 1$ and a length $n \ge 0$, and let $s \in (\bbz/k\bbz)^n$. Writing
$0$ for the zero vector of $(\bbz/k\bbz)^n$, the dot product of $s$ with $0$ equals $0$
in $\bbz/k\bbz$; that is, $\sum_{i=0}^{n-1} s_i \cdot 0 = 0$.
\end{lemma}

\begin{lemma}[Negating the left argument of the dot product]
\lean{ZkFourier.zkDot_neg_left}
\label{ZkFourier.zkDot_neg_left}
\leanok
\uses{ZkFourier.ZkVec, ZkFourier.zkDot}
\difficulty{2}
Fix an integer $k \ge 1$ and a dimension $n$, and work in the module $\mathbb{Z}_k^n =
(\mathbb{Z}/k\mathbb{Z})^n$ with dot product $s \cdot x = \sum_{i=0}^{n-1} s_i x_i \in
\mathbb{Z}/k\mathbb{Z}$. Then for all $s, x \in \mathbb{Z}_k^n$,
\[
(-s) \cdot x = -(s \cdot x).
\]
\end{lemma}

\begin{lemma}[Left additivity of the dot product under subtraction]
\lean{ZkFourier.zkDot_sub}
\label{ZkFourier.zkDot_sub}
\leanok
\uses{ZkFourier.ZkVec, ZkFourier.zkDot}
\difficulty{2}
For any three vectors $s, t, x \in \mathbb{Z}_k^n$, the dot product is additive in its
left argument under subtraction:
\[
(s - t) \cdot x = s \cdot x - t \cdot x,
\]
where the subtraction on the left is taken componentwise and both sides are elements of
$\mathbb{Z}/k\mathbb{Z}$.
\end{lemma}

\begin{lemma}[Additivity of a Fourier character on $\mathbb{Z}_k^n$]
\lean{ZkFourier.char_s_add}
\statementsource{boolean-ch08-generalized-domains}{
  pdf-50833afe7403-p021-b004,
  pdf-50833afe7403-p022-b009
}
\label{ZkFourier.char_s_add}
\leanok
\uses{ZkFourier.ZkVec, ZkFourier.char_s, ZkFourier.toOmega_add, ZkFourier.zkDot_add_right}
\difficulty{1}
Fix integers $k \ge 1$ and $n \ge 0$, let $\omega_k = e^{2\pi i/k}$, and for $s \in
(\mathbb{Z}/k\mathbb{Z})^n$ let $\chi_s : (\mathbb{Z}/k\mathbb{Z})^n \to \mathbb{C}$ be
the character $\chi_s(x) = \omega_k^{\,s\cdot x}$, where $s\cdot x = \sum_{i=1}^{n} s_i
x_i \in \mathbb{Z}/k\mathbb{Z}$ is the componentwise dot product. Then for all $x, y \in
(\mathbb{Z}/k\mathbb{Z})^n$,
\[
\chi_s(x + y) = \chi_s(x)\,\chi_s(y).
\]
\end{lemma}

\begin{lemma}[Character evaluated at the zero vector]
\lean{ZkFourier.char_s_zero_vec}
\label{ZkFourier.char_s_zero_vec}
\leanok
\uses{ZkFourier.ZkVec, ZkFourier.char_s, ZkFourier.toOmega_zero, ZkFourier.zkDot_zero_right}
\difficulty{1}
Fix an integer $k \ge 1$ and let $\omega_k = e^{2\pi i/k}$. For any $s \in
\mathbb{Z}_k^n$, the Fourier character $\chi_s(x) = \omega_k^{s \cdot x}$ satisfies
$\chi_s(0) = 1$, where $0$ denotes the zero vector of $\mathbb{Z}_k^n$.
\end{lemma}

\begin{lemma}[The zero-index character is identically one]
\lean{ZkFourier.char_s_zero_index}
\label{ZkFourier.char_s_zero_index}
\leanok
\uses{ZkFourier.ZkVec, ZkFourier.char_s, ZkFourier.toOmega_zero, ZkFourier.zkDot_zero_left}
\difficulty{1}
Fix an integer $k \ge 1$ and a dimension $n$, and let $\chi_s(x) = \omega_k^{s \cdot x}$
denote the Fourier character on $\mathbb{Z}_k^n$ indexed by $s \in \mathbb{Z}_k^n$,
where $\omega_k = e^{2\pi i/k}$ and $s \cdot x = \sum_{i} s_i x_i$ is the dot product
modulo $k$. Then the character indexed by the zero vector is the constant function $1$:
for every $x \in \mathbb{Z}_k^n$, \[ \chi_0(x) = 1. \]
\end{lemma}

\begin{lemma}[Fourier characters have unit modulus]
\lean{ZkFourier.norm_char_s}
\label{ZkFourier.norm_char_s}
\leanok
\uses{ZkFourier.ZkVec, ZkFourier.char_s, ZkFourier.norm_toOmega}
\difficulty{1}
Fix an integer $k \ge 1$ and let $\omega_k = e^{2\pi i/k}$ be the primitive $k$-th root
of unity. For $s \in \mathbb{Z}_k^n$, let $\chi_s \colon \mathbb{Z}_k^n \to \bbc$ be the
Fourier character $\chi_s(x) = \omega_k^{\,s\cdot x}$, where $s \cdot x = \sum_{i} s_i
x_i \in \mathbb{Z}/k\mathbb{Z}$ is the componentwise dot product. Then for every $s, x
\in \mathbb{Z}_k^n$ one has $\abs{\chi_s(x)} = 1$.
\end{lemma}

\begin{lemma}[Nonvanishing of the Fourier characters]
\lean{ZkFourier.char_s_ne_zero}
\label{ZkFourier.char_s_ne_zero}
\leanok
\uses{ZkFourier.ZkVec, ZkFourier.char_s, ZkFourier.norm_char_s}
\difficulty{2}
Fix $k \geq 1$ and $n \in \bbn$, and let $s \in \bbz_k^n$. Then the Fourier character
$\chi_s(x) = \omega_k^{s \cdot x}$ is nonzero for every $x \in \bbz_k^n$.
\end{lemma}

\begin{lemma}[Conjugate of a Fourier character]
\lean{ZkFourier.char_s_conj}
\statementsource{boolean-ch08-generalized-domains}{pdf-50833afe7403-p023-b002}
\label{ZkFourier.char_s_conj}
\leanok
\uses{ZkFourier.ZkVec, ZkFourier.char_s, ZkFourier.toOmega_neg, ZkFourier.zkDot_neg_left}
\difficulty{1}
Fix an integer $k \ge 1$, and for $s \in \mathbb{Z}_k^n$ let $\chi_s(x) = \omega_k^{s
\cdot x}$ be the Fourier character on $\mathbb{Z}_k^n$ indexed by $s$, where $\omega_k =
e^{2\pi i/k}$ and $s \cdot x = \sum_{i} s_i x_i \in \mathbb{Z}/k\mathbb{Z}$. Then for
every $s, x \in \mathbb{Z}_k^n$,
\[
\overline{\chi_s(x)} = \chi_{-s}(x);
\]
that is, complex conjugation of a character amounts to negating its index.
\end{lemma}

\begin{lemma}[Multiplicativity of Fourier characters]
\lean{ZkFourier.char_s_mul}
\proofsource{boolean-ch08-generalized-domains}{pdf-50833afe7403-p023-b002}
\label{ZkFourier.char_s_mul}
\leanok
\proofstep
  {ZkFourier.ZkVec}
  {context}
  {boolean-ch08-generalized-domains}
  {pdf-50833afe7403-p022-b009,
    pdf-50833afe7403-p021-b001}
\proofstep
  {ZkFourier.instDecidableEqZkVec}
  {context}
  {boolean-ch08-generalized-domains}
  {pdf-50833afe7403-p022-b009}
\proofstep
  {ZkFourier.instFintypeZkVecOfNeZeroNat}
  {context}
  {boolean-ch08-generalized-domains}
  {pdf-50833afe7403-p021-b002,
    pdf-50833afe7403-p022-b009}
\proofstep
  {ZkFourier.instAddCommGroupZkVec}
  {context}
  {boolean-ch08-generalized-domains}
  {pdf-50833afe7403-p021-b001,
    pdf-50833afe7403-p021-b002}
\proofstep
  {ZkFourier.instModuleZModZkVec}
  {context}
  {boolean-ch08-generalized-domains}
  {pdf-50833afe7403-p023-b007}
\proofstep
  {ZkFourier.rootOfUnity}
  {direct}
  {boolean-ch08-generalized-domains}
  {pdf-50833afe7403-p022-b007}
\proofstep
  {ZkFourier.toOmega}
  {direct}
  {boolean-ch08-generalized-domains}
  {pdf-50833afe7403-p022-b007}
\proofstep
  {ZkFourier.isPrimitiveRoot_rootOfUnity}
  {context}
  {boolean-ch08-generalized-domains}
  {pdf-50833afe7403-p022-b007}
\proofstep
  {ZkFourier.rootOfUnity_pow_k}
  {direct}
  {boolean-ch08-generalized-domains}
  {pdf-50833afe7403-p023-b001,
    pdf-50833afe7403-p021-b004}
\proofstep
  {ZkFourier.toOmega_zero}
  {direct}
  {boolean-ch08-generalized-domains}
  {pdf-50833afe7403-p021-b004}
\proofstep
  {ZkFourier.toOmega_add}
  {direct}
  {boolean-ch08-generalized-domains}
  {pdf-50833afe7403-p021-b004}
\proofstep
  {ZkFourier.ZkFun}
  {context}
  {boolean-ch08-generalized-domains}
  {pdf-50833afe7403-p021-b002}
\proofstep
  {ZkFourier.zkDot}
  {context}
  {boolean-ch08-generalized-domains}
  {pdf-50833afe7403-p022-b009}
\proofstep
  {ZkFourier.char_s}
  {direct}
  {boolean-ch08-generalized-domains}
  {pdf-50833afe7403-p022-b009,
    pdf-50833afe7403-p022-b007}
\proofstep
  {ZkFourier.zkDot_zero_left}
  {context}
  {boolean-ch08-generalized-domains}
  {pdf-50833afe7403-p022-b009}
\proofstep
  {ZkFourier.zkDot_zero_right}
  {context}
  {boolean-ch08-generalized-domains}
  {pdf-50833afe7403-p022-b009}
\proofstep
  {ZkFourier.char_s_zero_vec}
  {context}
  {boolean-ch08-generalized-domains}
  {pdf-50833afe7403-p021-b004}
\proofstep
  {ZkFourier.char_s_zero_index}
  {context}
  {boolean-ch08-generalized-domains}
  {pdf-50833afe7403-p022-b008}
\uses{ZkFourier.ZkVec, ZkFourier.char_s, ZkFourier.toOmega, ZkFourier.toOmega_add, ZkFourier.zkDot}
\difficulty{3}
Fix an integer $k \ge 1$ and let $\omega_k = e^{2\pi i/k}$. For a frequency vector $s
\in \mathbb{Z}_k^n$, let $\chi_s \colon \mathbb{Z}_k^n \to \mathbb{C}$ be the Fourier
character $\chi_s(x) = \omega_k^{\,s \cdot x}$, where $s \cdot x = \sum_{i} s_i x_i$ is
the dot product in $\mathbb{Z}/k\mathbb{Z}$. Then for all $s, t, x \in \mathbb{Z}_k^n$,
\[
\chi_s(x)\,\chi_t(x) = \chi_{s+t}(x),
\]
so that pointwise multiplication of characters corresponds to addition of their
frequency indices.
\end{lemma}

\begin{lemma}[Mean of a nontrivial Fourier character vanishes]
\lean{ZkFourier.expectation_char_nontrivial}
\statementsource{boolean-ch08-generalized-domains}{
  pdf-50833afe7403-p022-b002,
  pdf-50833afe7403-p022-b003
}
\label{ZkFourier.expectation_char_nontrivial}
\leanok
\uses{ZkFourier.ZkVec, ZkFourier.char_s, ZkFourier.expectation, ZkFourier.geom_sum_toOmega, ZkFourier.toOmega, ZkFourier.toOmega_add, ZkFourier.zkDot}
\difficulty{5}
Fix a modulus $k \ge 1$ and a dimension $n$, and let $\omega_k = e^{2\pi i/k}$. For a
frequency $s \in \bbz_k^n$, let $\chi_s(x) = \omega_k^{s\cdot x}$ be the associated
Fourier character on $\bbz_k^n$. If $s \neq 0$, then the character has vanishing mean:
\[
  \E[\chi_s] = \frac{1}{k^n}\sum_{x \in \bbz_k^n} \omega_k^{s\cdot x} = 0.
\]
\end{lemma}

\begin{lemma}[Expectation of the trivial character equals one]
\lean{ZkFourier.expectation_char_trivial}
\label{ZkFourier.expectation_char_trivial}
\leanok
\uses{ZkFourier.ZkVec, ZkFourier.card_ZkVec, ZkFourier.char_s, ZkFourier.expectation}
\difficulty{2}
Fix an integer $k \ge 1$ and let $\omega_k = e^{2\pi i/k}$ be the associated primitive
$k$-th root of unity. For a natural number $n$ and an index $s \in
(\mathbb{Z}/k\mathbb{Z})^n$, let $\chi_s \colon (\mathbb{Z}/k\mathbb{Z})^n \to
\mathbb{C}$ be the Fourier character $\chi_s(x) = \omega_k^{\,s \cdot x}$, and write
$\mathbb{E}[f] = k^{-n}\sum_{x} f(x)$ for the uniform average of a complex-valued
function $f$ over $(\mathbb{Z}/k\mathbb{Z})^n$. Then the trivial character $\chi_0$,
indexed by the zero vector, has expectation $\mathbb{E}[\chi_0] = 1$.
\end{lemma}

\begin{lemma}[Self-inner-product of a Fourier character]
\lean{ZkFourier.inner_product_char_self}
\proofsource{boolean-ch08-generalized-domains}{
  pdf-50833afe7403-p022-b004,
  pdf-50833afe7403-p022-b005
}
\label{ZkFourier.inner_product_char_self}
\leanok
\proofstep
  {ZkFourier.ZkVec}
  {context}
  {boolean-ch08-generalized-domains}
  {pdf-50833afe7403-p021-b002,
    pdf-50833afe7403-p022-b009}
\proofstep
  {ZkFourier.instDecidableEqZkVec}
  {context}
  {boolean-ch08-generalized-domains}
  {pdf-50833afe7403-p021-b002}
\proofstep
  {ZkFourier.instFintypeZkVecOfNeZeroNat}
  {context}
  {boolean-ch08-generalized-domains}
  {pdf-50833afe7403-p021-b002,
    pdf-50833afe7403-p021-b004}
\proofstep
  {ZkFourier.instAddCommGroupZkVec}
  {context}
  {boolean-ch08-generalized-domains}
  {pdf-50833afe7403-p021-b001}
\proofstep
  {ZkFourier.instModuleZModZkVec}
  {context}
  {boolean-ch08-generalized-domains}
  {pdf-50833afe7403-p023-b007}
\proofstep
  {ZkFourier.card_ZkVec}
  {context}
  {boolean-ch08-generalized-domains}
  {pdf-50833afe7403-p002-b006}
\proofstep
  {ZkFourier.rootOfUnity}
  {direct}
  {boolean-ch08-generalized-domains}
  {pdf-50833afe7403-p022-b007}
\proofstep
  {ZkFourier.toOmega}
  {context}
  {boolean-ch08-generalized-domains}
  {pdf-50833afe7403-p022-b007}
\proofstep
  {ZkFourier.toOmega_zero}
  {context}
  {boolean-ch08-generalized-domains}
  {pdf-50833afe7403-p021-b004}
\proofstep
  {ZkFourier.norm_toOmega}
  {direct}
  {boolean-ch08-generalized-domains}
  {pdf-50833afe7403-p021-b004}
\proofstep
  {ZkFourier.ZkFun}
  {context}
  {boolean-ch08-generalized-domains}
  {pdf-50833afe7403-p021-b002}
\proofstep
  {ZkFourier.expectation}
  {context}
  {boolean-ch08-generalized-domains}
  {pdf-50833afe7403-p021-b002}
\proofstep
  {ZkFourier.inner_product}
  {direct}
  {boolean-ch08-generalized-domains}
  {pdf-50833afe7403-p021-b002}
\proofstep
  {ZkFourier.zkDot}
  {context}
  {boolean-ch08-generalized-domains}
  {pdf-50833afe7403-p022-b009}
\proofstep
  {ZkFourier.char_s}
  {direct}
  {boolean-ch08-generalized-domains}
  {pdf-50833afe7403-p022-b009,
    pdf-50833afe7403-p022-b007}
\proofstep
  {ZkFourier.zkDot_zero_left}
  {context}
  {boolean-ch08-generalized-domains}
  {pdf-50833afe7403-p022-b009}
\proofstep
  {ZkFourier.zkDot_zero_right}
  {context}
  {boolean-ch08-generalized-domains}
  {pdf-50833afe7403-p022-b009}
\proofstep
  {ZkFourier.char_s_zero_vec}
  {context}
  {boolean-ch08-generalized-domains}
  {pdf-50833afe7403-p021-b004}
\proofstep
  {ZkFourier.char_s_zero_index}
  {context}
  {boolean-ch08-generalized-domains}
  {pdf-50833afe7403-p022-b008}
\proofstep
  {ZkFourier.norm_char_s}
  {direct}
  {boolean-ch08-generalized-domains}
  {pdf-50833afe7403-p021-b004}
\uses{ZkFourier.ZkVec, ZkFourier.card_ZkVec, ZkFourier.char_s, ZkFourier.expectation, ZkFourier.inner_product, ZkFourier.norm_char_s}
\difficulty{2}
Fix an integer $k \ge 1$ and a dimension $n \ge 0$, and let $\omega_k = e^{2\pi i/k}$.
For $s \in \mathbb{Z}_k^n$, let $\chi_s(x) = \omega_k^{\,s\cdot x}$ be the Fourier
character indexed by $s$, and equip functions on $\mathbb{Z}_k^n$ with the Hermitian
inner product $\langle f,g\rangle = \frac{1}{k^n}\sum_{x\in\mathbb{Z}_k^n}
f(x)\,\overline{g(x)}$. Then every character has unit norm: for all $s \in
\mathbb{Z}_k^n$, \[ \langle \chi_s, \chi_s\rangle = 1. \]
\end{lemma}

\begin{lemma}[Orthogonality of distinct characters]
\lean{ZkFourier.inner_product_char_nonself}
\proofsource{boolean-ch08-generalized-domains}{
  pdf-50833afe7403-p022-b004,
  pdf-50833afe7403-p022-b005
}
\label{ZkFourier.inner_product_char_nonself}
\leanok
\proofstep
  {ZkFourier.ZkVec}
  {context}
  {boolean-ch08-generalized-domains}
  {pdf-50833afe7403-p021-b002,
    pdf-50833afe7403-p022-b009}
\proofstep
  {ZkFourier.instDecidableEqZkVec}
  {context}
  {boolean-ch08-generalized-domains}
  {pdf-50833afe7403-p021-b002}
\proofstep
  {ZkFourier.instFintypeZkVecOfNeZeroNat}
  {context}
  {boolean-ch08-generalized-domains}
  {pdf-50833afe7403-p021-b002}
\proofstep
  {ZkFourier.instAddCommGroupZkVec}
  {context}
  {boolean-ch08-generalized-domains}
  {pdf-50833afe7403-p021-b001}
\proofstep
  {ZkFourier.instModuleZModZkVec}
  {context}
  {boolean-ch08-generalized-domains}
  {pdf-50833afe7403-p023-b007}
\proofstep
  {ZkFourier.rootOfUnity}
  {direct}
  {boolean-ch08-generalized-domains}
  {pdf-50833afe7403-p022-b007}
\proofstep
  {ZkFourier.toOmega}
  {context}
  {boolean-ch08-generalized-domains}
  {pdf-50833afe7403-p022-b007}
\proofstep
  {ZkFourier.isPrimitiveRoot_rootOfUnity}
  {context}
  {boolean-ch08-generalized-domains}
  {pdf-50833afe7403-p022-b007,
    pdf-50833afe7403-p022-b006}
\proofstep
  {ZkFourier.rootOfUnity_pow_k}
  {direct}
  {boolean-ch08-generalized-domains}
  {pdf-50833afe7403-p023-b001,
    pdf-50833afe7403-p022-b002}
\proofstep
  {ZkFourier.toOmega_zero}
  {context}
  {boolean-ch08-generalized-domains}
  {pdf-50833afe7403-p022-b002}
\proofstep
  {ZkFourier.toOmega_add}
  {direct}
  {boolean-ch08-generalized-domains}
  {pdf-50833afe7403-p021-b004}
\proofstep
  {ZkFourier.toOmega_neg}
  {context}
  {boolean-ch08-generalized-domains}
  {pdf-50833afe7403-p022-b005}
\proofstep
  {ZkFourier.toOmega_mul}
  {context}
  {boolean-ch08-generalized-domains}
  {pdf-50833afe7403-p022-b007}
\proofstep
  {ZkFourier.geom_sum_toOmega}
  {direct}
  {boolean-ch08-generalized-domains}
  {pdf-50833afe7403-p022-b002,
    pdf-50833afe7403-p022-b003}
\proofstep
  {ZkFourier.ZkFun}
  {context}
  {boolean-ch08-generalized-domains}
  {pdf-50833afe7403-p021-b002}
\proofstep
  {ZkFourier.expectation}
  {context}
  {boolean-ch08-generalized-domains}
  {pdf-50833afe7403-p021-b002}
\proofstep
  {ZkFourier.inner_product}
  {direct}
  {boolean-ch08-generalized-domains}
  {pdf-50833afe7403-p021-b002}
\proofstep
  {ZkFourier.zkDot}
  {context}
  {boolean-ch08-generalized-domains}
  {pdf-50833afe7403-p022-b009}
\proofstep
  {ZkFourier.char_s}
  {direct}
  {boolean-ch08-generalized-domains}
  {pdf-50833afe7403-p022-b009}
\proofstep
  {ZkFourier.zkDot_zero_left}
  {context}
  {boolean-ch08-generalized-domains}
  {pdf-50833afe7403-p022-b009}
\proofstep
  {ZkFourier.zkDot_zero_right}
  {context}
  {boolean-ch08-generalized-domains}
  {pdf-50833afe7403-p022-b002}
\proofstep
  {ZkFourier.zkDot_neg_left}
  {context}
  {boolean-ch08-generalized-domains}
  {pdf-50833afe7403-p023-b002}
\proofstep
  {ZkFourier.char_s_zero_vec}
  {context}
  {boolean-ch08-generalized-domains}
  {pdf-50833afe7403-p022-b002}
\proofstep
  {ZkFourier.char_s_zero_index}
  {context}
  {boolean-ch08-generalized-domains}
  {pdf-50833afe7403-p022-b008}
\proofstep
  {ZkFourier.char_s_conj}
  {context}
  {boolean-ch08-generalized-domains}
  {pdf-50833afe7403-p023-b002}
\proofstep
  {ZkFourier.char_s_mul}
  {direct}
  {boolean-ch08-generalized-domains}
  {pdf-50833afe7403-p023-b002}
\proofstep
  {ZkFourier.expectation_char_nontrivial}
  {direct}
  {boolean-ch08-generalized-domains}
  {pdf-50833afe7403-p022-b002,
    pdf-50833afe7403-p022-b003}
\uses{ZkFourier.ZkVec, ZkFourier.char_s, ZkFourier.char_s_conj, ZkFourier.char_s_mul, ZkFourier.expectation_char_nontrivial, ZkFourier.inner_product}
\difficulty{3}
Let $k \ge 1$ and $n \ge 0$, and let $\omega_k = e^{2\pi i/k}$. For $s, t \in
\mathbb{Z}_k^n$ with $s \ne t$, the Fourier characters $\chi_s(x) = \omega_k^{s\cdot x}$
and $\chi_t(x) = \omega_k^{t\cdot x}$ are orthogonal with respect to the Hermitian inner
product $\langle f, g\rangle = \frac{1}{k^n}\sum_{x \in \mathbb{Z}_k^n}
f(x)\,\overline{g(x)}$; that is, $\langle \chi_s, \chi_t\rangle = 0$.
\end{lemma}

\begin{definition}[Fourier coefficient]
\lean{ZkFourier.fourier_coeff}
\label{ZkFourier.fourier_coeff}
\leanok
\statementsource{boolean-ch08-generalized-domains}{pdf-50833afe7403-p005-b001}
\uses{ZkFourier.ZkFun, ZkFourier.ZkVec, ZkFourier.char_s, ZkFourier.inner_product}
The Fourier coefficient of $f : \mathbb{Z}_k^n \to \mathbb{C}$ at frequency
$s$ is
\[
  \hat{f}(s) = \langle f, \chi_s \rangle
             = \mathbb{E}_{x}\!\left[f(x)\,\overline{\chi_s(x)}\right].
\]
\end{definition}

\begin{lemma}[Character has unit Fourier coefficient at its own frequency]
\lean{ZkFourier.fourier_coeff_char_self}
\proofsource{boolean-ch08-generalized-domains}{
  pdf-50833afe7403-p022-b005,
  pdf-50833afe7403-p022-b004
}
\label{ZkFourier.fourier_coeff_char_self}
\leanok
\proofstep
  {ZkFourier.ZkVec}
  {context}
  {boolean-ch08-generalized-domains}
  {pdf-50833afe7403-p021-b002,
    pdf-50833afe7403-p021-b001}
\proofstep
  {ZkFourier.instDecidableEqZkVec}
  {context}
  {boolean-ch08-generalized-domains}
  {pdf-50833afe7403-p021-b002}
\proofstep
  {ZkFourier.instFintypeZkVecOfNeZeroNat}
  {context}
  {boolean-ch08-generalized-domains}
  {pdf-50833afe7403-p021-b002}
\proofstep
  {ZkFourier.instAddCommGroupZkVec}
  {context}
  {boolean-ch08-generalized-domains}
  {pdf-50833afe7403-p021-b001}
\proofstep
  {ZkFourier.instModuleZModZkVec}
  {context}
  {boolean-ch08-generalized-domains}
  {pdf-50833afe7403-p023-b007}
\proofstep
  {ZkFourier.card_ZkVec}
  {context}
  {boolean-ch08-generalized-domains}
  {pdf-50833afe7403-p002-b006}
\proofstep
  {ZkFourier.rootOfUnity}
  {direct}
  {boolean-ch08-generalized-domains}
  {pdf-50833afe7403-p022-b007}
\proofstep
  {ZkFourier.toOmega}
  {direct}
  {boolean-ch08-generalized-domains}
  {pdf-50833afe7403-p022-b007}
\proofstep
  {ZkFourier.toOmega_zero}
  {context}
  {boolean-ch08-generalized-domains}
  {pdf-50833afe7403-p021-b004}
\proofstep
  {ZkFourier.norm_toOmega}
  {direct}
  {boolean-ch08-generalized-domains}
  {pdf-50833afe7403-p021-b004}
\proofstep
  {ZkFourier.ZkFun}
  {context}
  {boolean-ch08-generalized-domains}
  {pdf-50833afe7403-p021-b002}
\proofstep
  {ZkFourier.expectation}
  {context}
  {boolean-ch08-generalized-domains}
  {pdf-50833afe7403-p021-b002}
\proofstep
  {ZkFourier.inner_product}
  {direct}
  {boolean-ch08-generalized-domains}
  {pdf-50833afe7403-p021-b002}
\proofstep
  {ZkFourier.zkDot}
  {context}
  {boolean-ch08-generalized-domains}
  {pdf-50833afe7403-p022-b009}
\proofstep
  {ZkFourier.char_s}
  {direct}
  {boolean-ch08-generalized-domains}
  {pdf-50833afe7403-p022-b009,
    pdf-50833afe7403-p022-b007}
\proofstep
  {ZkFourier.zkDot_zero_left}
  {context}
  {boolean-ch08-generalized-domains}
  {pdf-50833afe7403-p022-b009}
\proofstep
  {ZkFourier.zkDot_zero_right}
  {context}
  {boolean-ch08-generalized-domains}
  {pdf-50833afe7403-p022-b009}
\proofstep
  {ZkFourier.char_s_zero_vec}
  {context}
  {boolean-ch08-generalized-domains}
  {pdf-50833afe7403-p021-b004}
\proofstep
  {ZkFourier.char_s_zero_index}
  {context}
  {boolean-ch08-generalized-domains}
  {pdf-50833afe7403-p022-b008}
\proofstep
  {ZkFourier.norm_char_s}
  {direct}
  {boolean-ch08-generalized-domains}
  {pdf-50833afe7403-p021-b004}
\proofstep
  {ZkFourier.inner_product_char_self}
  {direct}
  {boolean-ch08-generalized-domains}
  {pdf-50833afe7403-p022-b005,
    pdf-50833afe7403-p022-b004}
\proofstep
  {ZkFourier.fourier_coeff}
  {direct}
  {boolean-ch08-generalized-domains}
  {pdf-50833afe7403-p005-b001}
\uses{ZkFourier.ZkVec, ZkFourier.char_s, ZkFourier.fourier_coeff, ZkFourier.inner_product_char_self}
\difficulty{1}
Fix an integer $k \ge 1$ and a dimension $n$, and let $\chi_s(x) = \omega_k^{s \cdot x}$
denote the Fourier character on $\mathbb{Z}_k^n$ indexed by a frequency $s \in
\mathbb{Z}_k^n$, where $\omega_k = e^{2\pi i / k}$. Then for every $s \in
\mathbb{Z}_k^n$, the Fourier coefficient of $\chi_s$ at its own frequency equals $1$;
equivalently, $\widehat{\chi_s}(s) = \langle \chi_s, \chi_s \rangle = 1$.
\end{lemma}

\begin{lemma}[Orthogonality of distinct Fourier characters]
\lean{ZkFourier.fourier_coeff_char_of_ne}
\proofsource{boolean-ch08-generalized-domains}{
  pdf-50833afe7403-p022-b004,
  pdf-50833afe7403-p022-b005
}
\label{ZkFourier.fourier_coeff_char_of_ne}
\leanok
\proofstep
  {ZkFourier.ZkVec}
  {context}
  {boolean-ch08-generalized-domains}
  {pdf-50833afe7403-p022-b009,
    pdf-50833afe7403-p021-b002}
\proofstep
  {ZkFourier.instDecidableEqZkVec}
  {context}
  {boolean-ch08-generalized-domains}
  {pdf-50833afe7403-p022-b009}
\proofstep
  {ZkFourier.instFintypeZkVecOfNeZeroNat}
  {context}
  {boolean-ch08-generalized-domains}
  {pdf-50833afe7403-p021-b002,
    pdf-50833afe7403-p021-b004}
\proofstep
  {ZkFourier.instAddCommGroupZkVec}
  {context}
  {boolean-ch08-generalized-domains}
  {pdf-50833afe7403-p021-b001,
    pdf-50833afe7403-p022-b009}
\proofstep
  {ZkFourier.instModuleZModZkVec}
  {context}
  {boolean-ch08-generalized-domains}
  {pdf-50833afe7403-p023-b007}
\proofstep
  {ZkFourier.rootOfUnity}
  {direct}
  {boolean-ch08-generalized-domains}
  {pdf-50833afe7403-p022-b007}
\proofstep
  {ZkFourier.toOmega}
  {context}
  {boolean-ch08-generalized-domains}
  {pdf-50833afe7403-p022-b007,
    pdf-50833afe7403-p022-b004}
\proofstep
  {ZkFourier.isPrimitiveRoot_rootOfUnity}
  {context}
  {boolean-ch08-generalized-domains}
  {pdf-50833afe7403-p022-b007,
    pdf-50833afe7403-p022-b004}
\proofstep
  {ZkFourier.rootOfUnity_pow_k}
  {direct}
  {boolean-ch08-generalized-domains}
  {pdf-50833afe7403-p023-b001,
    pdf-50833afe7403-p022-b004}
\proofstep
  {ZkFourier.toOmega_zero}
  {context}
  {boolean-ch08-generalized-domains}
  {pdf-50833afe7403-p022-b004}
\proofstep
  {ZkFourier.toOmega_add}
  {direct}
  {boolean-ch08-generalized-domains}
  {pdf-50833afe7403-p022-b004}
\proofstep
  {ZkFourier.toOmega_neg}
  {context}
  {boolean-ch08-generalized-domains}
  {pdf-50833afe7403-p023-b001,
    pdf-50833afe7403-p022-b001}
\proofstep
  {ZkFourier.toOmega_mul}
  {direct}
  {boolean-ch08-generalized-domains}
  {pdf-50833afe7403-p022-b007}
\proofstep
  {ZkFourier.geom_sum_toOmega}
  {direct}
  {boolean-ch08-generalized-domains}
  {pdf-50833afe7403-p022-b002,
    pdf-50833afe7403-p022-b003}
\proofstep
  {ZkFourier.ZkFun}
  {context}
  {boolean-ch08-generalized-domains}
  {pdf-50833afe7403-p021-b002}
\proofstep
  {ZkFourier.expectation}
  {context}
  {boolean-ch08-generalized-domains}
  {pdf-50833afe7403-p021-b002}
\proofstep
  {ZkFourier.inner_product}
  {direct}
  {boolean-ch08-generalized-domains}
  {pdf-50833afe7403-p021-b002}
\proofstep
  {ZkFourier.zkDot}
  {context}
  {boolean-ch08-generalized-domains}
  {pdf-50833afe7403-p022-b009}
\proofstep
  {ZkFourier.char_s}
  {direct}
  {boolean-ch08-generalized-domains}
  {pdf-50833afe7403-p022-b009,
    pdf-50833afe7403-p022-b007}
\proofstep
  {ZkFourier.zkDot_zero_left}
  {context}
  {boolean-ch08-generalized-domains}
  {pdf-50833afe7403-p022-b008}
\proofstep
  {ZkFourier.zkDot_zero_right}
  {context}
  {boolean-ch08-generalized-domains}
  {pdf-50833afe7403-p022-b004}
\proofstep
  {ZkFourier.zkDot_neg_left}
  {context}
  {boolean-ch08-generalized-domains}
  {pdf-50833afe7403-p023-b002}
\proofstep
  {ZkFourier.char_s_zero_vec}
  {context}
  {boolean-ch08-generalized-domains}
  {pdf-50833afe7403-p022-b004}
\proofstep
  {ZkFourier.char_s_zero_index}
  {context}
  {boolean-ch08-generalized-domains}
  {pdf-50833afe7403-p022-b008}
\proofstep
  {ZkFourier.char_s_conj}
  {direct}
  {boolean-ch08-generalized-domains}
  {pdf-50833afe7403-p023-b002,
    pdf-50833afe7403-p022-b001}
\proofstep
  {ZkFourier.char_s_mul}
  {direct}
  {boolean-ch08-generalized-domains}
  {pdf-50833afe7403-p023-b002}
\proofstep
  {ZkFourier.expectation_char_nontrivial}
  {direct}
  {boolean-ch08-generalized-domains}
  {pdf-50833afe7403-p022-b002,
    pdf-50833afe7403-p022-b003}
\proofstep
  {ZkFourier.inner_product_char_nonself}
  {direct}
  {boolean-ch08-generalized-domains}
  {pdf-50833afe7403-p022-b004,
    pdf-50833afe7403-p022-b005}
\proofstep
  {ZkFourier.fourier_coeff}
  {direct}
  {boolean-ch08-generalized-domains}
  {pdf-50833afe7403-p021-b003,
    pdf-50833afe7403-p005-b001}
\uses{ZkFourier.ZkVec, ZkFourier.char_s, ZkFourier.fourier_coeff, ZkFourier.inner_product_char_nonself}
\difficulty{1}
Fix an integer $k \ge 1$ and a dimension $n$, and let $\omega_k = e^{2\pi i / k}$. For a
frequency $s \in \mathbb{Z}_k^n$, write $\chi_s(x) = \omega_k^{\,s \cdot x}$ for the
associated Fourier character on $\mathbb{Z}_k^n$, where $s \cdot x = \sum_{i} s_i x_i$
is computed in $\mathbb{Z}/k\mathbb{Z}$. Then for any two distinct frequencies $s \ne t$
in $\mathbb{Z}_k^n$, the Fourier coefficient of the character $\chi_s$ at $t$ vanishes:
\[
\widehat{\chi_s}(t) = \langle \chi_s, \chi_t \rangle =
\mathbb{E}_{x}\!\left[\chi_s(x)\,\overline{\chi_t(x)}\right] = 0.
\]
\end{lemma}

\begin{lemma}[Fourier expansion on $\mathbb{Z}_k^n$]
\lean{ZkFourier.fourier_expansion}
\label{ZkFourier.fourier_expansion}
\leanok
\statementsource{boolean-ch08-generalized-domains}{
  pdf-50833afe7403-p005-b001,
  pdf-50833afe7403-p022-b009
}
\uses{ZkFourier.ZkFun, ZkFourier.ZkVec, ZkFourier.card_ZkVec, ZkFourier.char_s, ZkFourier.char_s_add, ZkFourier.char_s_conj, ZkFourier.expectation, ZkFourier.expectation_char_nontrivial, ZkFourier.fourier_coeff, ZkFourier.inner_product, ZkFourier.zkDot}
\difficulty{6}
Fix integers $k \ge 1$ and $n$, and let $\omega_k = e^{2\pi i/k}$. For a frequency $s
\in \mathbb{Z}_k^n$, let $\chi_s(x) = \omega_k^{\,s\cdot x}$ be the corresponding
Fourier character, where $s\cdot x = \sum_{i} s_i x_i \in \mathbb{Z}/k\mathbb{Z}$, and
let $\hat f(s) = \tfrac{1}{k^n}\sum_{x \in \mathbb{Z}_k^n} f(x)\,\overline{\chi_s(x)}$
be the Fourier coefficient of $f$ at $s$. Then every function $f : \mathbb{Z}_k^n \to
\mathbb{C}$ is recovered from its Fourier coefficients by
\[
  f(x) = \sum_{s \in \mathbb{Z}_k^n} \hat f(s)\,\chi_s(x)
  \qquad\text{for all } x \in \mathbb{Z}_k^n.
\]
\end{lemma}

\begin{lemma}[Parseval's identity on $\mathbb{Z}_k^n$]
\lean{ZkFourier.parseval_identity}
\proofsource{boolean-ch08-generalized-domains}{
  pdf-50833afe7403-p006-b001,
  pdf-50833afe7403-p005-b006,
  pdf-50833afe7403-p021-b003
}
\label{ZkFourier.parseval_identity}
\leanok
\proofstep
  {ZkFourier.ZkVec}
  {context}
  {boolean-ch08-generalized-domains}
  {pdf-50833afe7403-p021-b002,
    pdf-50833afe7403-p022-b009}
\proofstep
  {ZkFourier.instDecidableEqZkVec}
  {context}
  {boolean-ch08-generalized-domains}
  {pdf-50833afe7403-p021-b002}
\proofstep
  {ZkFourier.instFintypeZkVecOfNeZeroNat}
  {context}
  {boolean-ch08-generalized-domains}
  {pdf-50833afe7403-p021-b002}
\proofstep
  {ZkFourier.instAddCommGroupZkVec}
  {context}
  {boolean-ch08-generalized-domains}
  {pdf-50833afe7403-p021-b001}
\proofstep
  {ZkFourier.instModuleZModZkVec}
  {context}
  {boolean-ch08-generalized-domains}
  {pdf-50833afe7403-p023-b007}
\proofstep
  {ZkFourier.card_ZkVec}
  {context}
  {boolean-ch08-generalized-domains}
  {pdf-50833afe7403-p002-b006}
\proofstep
  {ZkFourier.rootOfUnity}
  {direct}
  {boolean-ch08-generalized-domains}
  {pdf-50833afe7403-p022-b007}
\proofstep
  {ZkFourier.toOmega}
  {context}
  {boolean-ch08-generalized-domains}
  {pdf-50833afe7403-p022-b007}
\proofstep
  {ZkFourier.isPrimitiveRoot_rootOfUnity}
  {context}
  {boolean-ch08-generalized-domains}
  {pdf-50833afe7403-p022-b007,
    pdf-50833afe7403-p022-b006}
\proofstep
  {ZkFourier.rootOfUnity_pow_k}
  {context}
  {boolean-ch08-generalized-domains}
  {pdf-50833afe7403-p023-b001,
    pdf-50833afe7403-p022-b002}
\proofstep
  {ZkFourier.toOmega_zero}
  {context}
  {boolean-ch08-generalized-domains}
  {pdf-50833afe7403-p022-b002}
\proofstep
  {ZkFourier.toOmega_add}
  {context}
  {boolean-ch08-generalized-domains}
  {pdf-50833afe7403-p022-b002}
\proofstep
  {ZkFourier.toOmega_neg}
  {context}
  {boolean-ch08-generalized-domains}
  {pdf-50833afe7403-p022-b001,
    pdf-50833afe7403-p023-b001}
\proofstep
  {ZkFourier.toOmega_mul}
  {context}
  {boolean-ch08-generalized-domains}
  {pdf-50833afe7403-p022-b007}
\proofstep
  {ZkFourier.geom_sum_toOmega}
  {direct}
  {boolean-ch08-generalized-domains}
  {pdf-50833afe7403-p022-b002,
    pdf-50833afe7403-p022-b003}
\proofstep
  {ZkFourier.ZkFun}
  {context}
  {boolean-ch08-generalized-domains}
  {pdf-50833afe7403-p021-b002}
\proofstep
  {ZkFourier.expectation}
  {context}
  {boolean-ch08-generalized-domains}
  {pdf-50833afe7403-p021-b002}
\proofstep
  {ZkFourier.inner_product}
  {direct}
  {boolean-ch08-generalized-domains}
  {pdf-50833afe7403-p021-b002}
\proofstep
  {ZkFourier.L2_norm_sq}
  {context}
  {boolean-ch08-generalized-domains}
  {pdf-50833afe7403-p021-b002}
\proofstep
  {ZkFourier.zkDot}
  {context}
  {boolean-ch08-generalized-domains}
  {pdf-50833afe7403-p022-b009}
\proofstep
  {ZkFourier.char_s}
  {direct}
  {boolean-ch08-generalized-domains}
  {pdf-50833afe7403-p022-b009,
    pdf-50833afe7403-p022-b007}
\proofstep
  {ZkFourier.zkDot_add_right}
  {context}
  {boolean-ch08-generalized-domains}
  {pdf-50833afe7403-p022-b009}
\proofstep
  {ZkFourier.zkDot_zero_left}
  {context}
  {boolean-ch08-generalized-domains}
  {pdf-50833afe7403-p022-b008}
\proofstep
  {ZkFourier.zkDot_zero_right}
  {context}
  {boolean-ch08-generalized-domains}
  {pdf-50833afe7403-p022-b009}
\proofstep
  {ZkFourier.zkDot_neg_left}
  {context}
  {boolean-ch08-generalized-domains}
  {pdf-50833afe7403-p023-b002}
\proofstep
  {ZkFourier.char_s_add}
  {context}
  {boolean-ch08-generalized-domains}
  {pdf-50833afe7403-p022-b002}
\proofstep
  {ZkFourier.char_s_zero_vec}
  {context}
  {boolean-ch08-generalized-domains}
  {pdf-50833afe7403-p022-b002}
\proofstep
  {ZkFourier.char_s_zero_index}
  {context}
  {boolean-ch08-generalized-domains}
  {pdf-50833afe7403-p022-b008}
\proofstep
  {ZkFourier.char_s_conj}
  {context}
  {boolean-ch08-generalized-domains}
  {pdf-50833afe7403-p023-b002}
\proofstep
  {ZkFourier.expectation_char_nontrivial}
  {direct}
  {boolean-ch08-generalized-domains}
  {pdf-50833afe7403-p022-b002,
    pdf-50833afe7403-p022-b003}
\proofstep
  {ZkFourier.fourier_coeff}
  {direct}
  {boolean-ch08-generalized-domains}
  {pdf-50833afe7403-p005-b001}
\proofstep
  {ZkFourier.fourier_expansion}
  {direct}
  {boolean-ch08-generalized-domains}
  {pdf-50833afe7403-p005-b001,
    pdf-50833afe7403-p022-b009}
\uses{ZkFourier.L2_norm_sq, ZkFourier.ZkFun, ZkFourier.ZkVec, ZkFourier.char_s, ZkFourier.fourier_coeff, ZkFourier.fourier_expansion}
\difficulty{5}
Let $k \ge 1$ and $n$ be natural numbers, and let $f : \mathbb{Z}_k^n \to \mathbb{C}$ be
a complex-valued function on $\mathbb{Z}_k^n$. Then the sum of the squared magnitudes of
its Fourier coefficients equals its squared $L^2$ norm:
\[
  \sum_{s \in \mathbb{Z}_k^n} \bigl|\hat{f}(s)\bigr|^2 = \|f\|_2^2,
\]
where $\hat{f}(s) = \mathbb{E}_{x}\!\left[f(x)\,\overline{\omega_k^{\,s \cdot
x}}\right]$ is the Fourier coefficient at frequency $s$ and $\|f\|_2^2 =
\tfrac{1}{k^n}\sum_{x \in \mathbb{Z}_k^n} |f(x)|^2$.
\end{lemma}

\begin{lemma}[Convolution theorem for the Fourier transform on $\mathbb{Z}_k^n$]
\lean{ZkFourier.fourier_coeff_convolution}
\statementsource{boolean-ch08-generalized-domains}{pdf-50833afe7403-p023-b006}
\label{ZkFourier.fourier_coeff_convolution}
\leanok
\uses{ZkFourier.ZkFun, ZkFourier.ZkVec, ZkFourier.char_s_add, ZkFourier.convolution, ZkFourier.expectation, ZkFourier.fourier_coeff, ZkFourier.inner_product}
\difficulty{4}
Fix an integer $k \ge 1$ and let $f, g : \mathbb{Z}_k^n \to \mathbb{C}$ be
complex-valued functions on $\mathbb{Z}_k^n$, with convolution $(f * g)(x) =
\mathbb{E}_y\!\left[f(y)\,g(x - y)\right]$ and Fourier coefficients $\hat{f}(s) =
\langle f, \chi_s\rangle$ taken at any frequency $s \in \mathbb{Z}_k^n$. Then the
Fourier transform carries convolution to pointwise multiplication:
\[
  \widehat{f * g}(s) = \hat{f}(s)\,\hat{g}(s).
\]
\end{lemma}
